\chapter{Stato dell'arte}
\label{chap:stato_arte}


Una delle domande più ricorrenti da parte degli studenti \`e cosa si intenda per ``stato dell'arte''. Concretamente, si tratta di individuare la situazione corrente riguardo la tematica trattata nell'elaborato. A titolo di esempio, se il lavoro di tesi si concentra su un (presunto) innovativo algoritmo per suggerire i brani di una playlist, è necessario verificare che tale approccio sia veramente innovativo e, in ogni caso, studiare gli approcci alternativi già disponibili in letteratura.

Per ottenere tale risultato, è necessario condurre ricerche approfondite a livello bibliografico (testi di riferimento, articoli scientifici, ecc.) e implementativo (sul web, sugli store, ecc.) \textbf{prima} di iniziare il proprio lavoro.

\section{Risorse}

Riguardo gli articoli scientifici, si caldeggia l'uso del motore di ricerca specializzato Google Scholar.\footnote{\url{https://scholar.google.it/}.} Esistono, poi, numerosi repository che consentono di accedere gratuitamente a pubblicazioni scientifiche, tra cui ResearchGate,\footnote{\url{https://www.researchgate.net/}} Zenodo,\footnote{\url{https://zenodo.org/}} e SBA - Sistema Bibliotecario di Ateneo.\footnote{\url{http://www.sba.unimi.it/index.html} con contenuti scaricabili solo dalla rete interna all'ateneo.} In particolare quest'ultimo permette agli studenti dell'Ateneo di accedere gratuitamente a molte pubblicazioni altrimenti accessibili solo a pagamento.

\`E importante riconoscere, nella grande mole di letteratura scientifica disponibile, i lavori pi\`u autorevoli e affidabili. Un primo elemento per orientarsi \`e che non tutte le sedi di pubblicazione sono uguali. In particolora una regola generale (per quanto non assoluta) \`e che gli articoli pubblicati su {\em rivista} sono tipicamente pi\`u completi e rigorosi di quelli pubblicati su atti di {\em convegno}. Un secondo elemento correlato all'autorevolezza della pubblicazione \`e il numero di citazioni che questa ha raccolto da parte di altri ricercatori (dato visibile ad esempio su Google Scholar). Infine, \`e opportuno riconoscere le sedi di pubblicazione (riviste e convegni) rilevanti per il settore: per questo, si consiglia di confrontarsi con i propri relatori e di consultare indicatori di qualità come Scimago JR\footnote{\url{https://www.scimagojr.com}} per le riviste e il GII-GRIN-SCIE Conference Rating\footnote{\url{https://scie.lcc.uma.es:8443}} per le conferenze.

\section{Buone pratiche}

Alcuni consigli di buone pratiche per l'analisi dello stato dell'arte.
\begin{itemize}
\item Imparare a leggere in maniera {\em efficace}. Ci sono migliaia di articoli potenzialmente interessanti, quindi \`e essenziale riuscire a estrarre in breve tempo gli elementi rilevanti di un articolo senza perdersi nei dettagli (fino a che non sia strettamente necessario), nonch\'e annotare in maniera ordinata tali elementi.
\item Imparare a leggere in maniera {\em critica}. Raramente un articolo scientifico \`e perfetto, gli aspetti tecnici e gli eventuali punti deboli (metodologie, ripetibilit\`a dei risultati, ecc.) vanno valutati e annotati con attenzione.
\item Imparare a seguire le ``piste'' interessanti. Una volta trovato un articolo rilevante, \`e utile esaminare
\begin{itemize}
\item gli articoli che esso cita (i suoi riferimenti bibliografici): questo permette di rintracciare riferimenti autorevoli a cui l'articolo si appoggia;
\item gli articoli che lo citano (ad esempio Google Scholar offre questa funzionalit\`a): questo permette di rintracciare riferimenti pi\`u recenti che hanno proseguito nella stessa direzione di ricerca.
\end{itemize}
\item Imparare a lavorare iterativamente.
\begin{itemize}
\item Aggiungere alla propria bibliografia gli articoli considerati rilevanti, a mano a mano che vengono trovati.
\item Raggrupparli iterativamente in sottotematiche.
\item Usare inizialmente un approccio ``inclusivo'' (nel dubbio, aggiungere un articolo in bibliografia piuttosto che scartarlo), e solo in un secondo tempo decidere cosa tenere e cosa scartare.
\end{itemize}
\end{itemize}



\section{Bibliografia e sitografia}
\label{sec:biblio}

La bibliografia di un lavoro scientifico deve necessariamente essere ricca. Tutti i testi in bibliografia devono essere citati almeno una volta nel corso dell'elaborato. 

Nella valutazione della bibliografia da parte della commissione, i testi cartacei sono considerati significativamente più validi dei siti consultati. Ne consegue che la bibliografia debba essere ricca di testi pubblicati, siano essi libri, articoli, al limite tesi di laurea o di dottorato o rapporti tecnici. Si suggerisce di evitare citazioni a fonti di dubbia valenza scientifica, quali Wikipedia, W3Schools, ecc.

Se è necessario citare siti Web, esistono tre strade ugualmente accettabili:
\begin{enumerate}
	\item se il numero di siti non è preponderante rispetto ai testi ``tradizionali'', è possibile inserirli parimenti in bibliografia;
	\item se i siti da citare sono numerosi, è più opportuno creare una sorta di bibliografia parallela e separata, detta \textit{sitografia};
	\item infine, se la citazione dei siti serve a individuare un prodotto e non una fonte di informazioni, la soluzione più opportuna è quella delle note a piè di pagina.
\end{enumerate}
